\documentclass[10pt,article]{article}
\usepackage[utf8]{inputenc}
\usepackage[margin=0.75in]{geometry}
\usepackage[fleqn]{amsmath}
\usepackage{amssymb}
\usepackage{booktabs}

\title{\vspace{-1.5cm}Holographic Storage Quick Reference}
\author{}
\date{}

\begin{document}

\maketitle
\vspace{-1cm}

\subsection*{Key Formulas}
\begin{flalign*}
    &\Delta\theta_{\text{FWHM}} \approx \frac{\lambda}{2 n L \cos\theta_B} &\\
    &\kappa = \frac{\pi \Delta n}{\lambda \cos\theta_B} &\\
    &\eta_{\max} = \sin^2(\kappa L) &\\
    &N_{\text{pages}} \approx \frac{\Theta_{\text{total}}}{2 \Delta\theta_{\text{FWHM}}} &
\end{flalign*}

\subsection*{Design Workflow}
\begin{enumerate}
    \setlength{\itemsep}{0pt}
    \setlength{\parskip}{0pt}
    \item Pick laser wavelength ($\lambda$).
    \item Choose material ($n$, $\Delta n$).
    \item Set thickness ($L$).
    \item Choose Bragg angle ($\theta_B$).
    \item Compute $\Delta\theta_{\text{FWHM}}$.
    \item Compute $\kappa$ and $\eta_{\max}$.
    \item Set angular budget $\Theta_{\text{total}}$.
    \item Estimate max pages ($N$).
    \item Adjust $L$, $\Delta n$, $\theta_B$ to balance capacity vs.\ efficiency.
\end{enumerate}

\subsection*{Common Materials}
\noindent
\begin{tabular}{@{} lccc @{}}
    \toprule
    \textbf{Material} & \textbf{$n$} & \textbf{$\Delta n$ range} & \textbf{Notes} \\
    \midrule
    Photopolymer & $1.5\text{--}1.6$ & $1\times 10^{-4} \text{--} 1\times 10^{-3}$ & High $\Delta n$, write-once \\
    $\text{LiNbO}_3$ (doped) & $\sim 2.2$ & $1\times 10^{-5} \text{--} 1\times 10^{-4}$ & Rewritable, low $\Delta n$ \\
    $\text{BaTiO}_3$ & $\sim 2.4$ & $1\times 10^{-5} \text{--} 1\times 10^{-4}$ & High sensitivity, unstable \\
    PMMA (doped) & $\sim 1.49$ & $1\times 10^{-5} \text{--} 1\times 10^{-4}$ & Stable, low sensitivity \\
    \bottomrule
\end{tabular}

\subsection*{Rules of Thumb}
\begin{itemize}
    \setlength{\itemsep}{0pt}
    \setlength{\parskip}{0pt}
    \item \textbf{Thicker} = more pages, sharper selectivity.
    \item \textbf{Higher $\Delta n$} = stronger efficiency.
    \item \textbf{Shorter $\lambda$} = narrower selectivity.
\end{itemize}

\end{document}